% !TeX root = GeoDynamics.tex
\newpage
\section{ShadowLinDynamics Source Flow}

\newcommand{\ShadowFunctionSection}[1]{%
  \section{ShadowLinDynamics:#1}%
  \label{shadowlin:#1}%
}

\texttt{ShadowLinDynamics} processes one source particle completely before
advancing to the next source.  It may read target frame-start state, but all
dynamics writes, contact slots, collision history, and internal-momentum
ledgers remain owned by the source particle.

\begin{figure}[ht]
\centering
\resizebox{\textwidth}{!}{\begin{tikzpicture}[
    node distance=0.34in and 0.62in,
    function/.style={
        draw=black!60,
        fill=black!3,
        rounded corners=2pt,
        minimum width=2.25in,
        minimum height=0.28in,
        align=center,
        font=\ttfamily\small
    },
    flow/.style={-{Stealth[length=5pt]}, draw=black!65, line width=0.45pt},
    loop/.style={flow, rounded corners=0pt}
]

\newcommand{\FlowNode}[2]{\hyperref[shadowlin:#1]{\strut #2}}

\node[function] (contact) {\FlowNode{ContactDynamics}{ContactDynamics}};
\node[function, below=of contact] (begin) {\FlowNode{BeginSourceContactParms}{BeginSourceContactParms}};
\node[function, below=of begin] (init) {\FlowNode{InitializeContactState}{InitializeContactState}};
\node[function, below=of init] (state) {\FlowNode{GetContactState}{GetContactState}};
\node[function, below=of state] (append) {\FlowNode{AppendContactSlot}{AppendContactSlot}};
\node[function, below=of append] (slots) {\FlowNode{GetContactSlots}{GetContactSlots}};
\node[function, below=of slots] (phase) {\FlowNode{GetCollisionPhase}{GetCollisionPhase}};
\node[function, below=of phase] (collisionmom) {\FlowNode{GetCollisionInternalMomentum}{GetCollisionInternalMomentum}};
\node[function, below=of collisionmom] (geometry) {\FlowNode{GetParticleGeometry}{GetParticleGeometry}};
\node[function, below=of geometry] (position) {\FlowNode{GetParticlePosition}{GetParticlePosition}};
\node[function, below=of position] (startvelocity) {\FlowNode{GetStartFrameVelocity}{GetStartFrameVelocity}};
\node[function, below=of startvelocity] (reducedmass) {\FlowNode{GetReducedMass}{GetReducedMass}};
\node[function, below=of reducedmass] (parms) {\FlowNode{CalculateContactParms}{CalculateContactParms}};

\node[function, right=of contact] (weights) {\FlowNode{CalculateContactWeights}{CalculateContactWeights}};
\node[function, below=of weights] (impulses) {\FlowNode{CalculateContactImpulses}{CalculateContactImpulses}};
\node[function, below=of impulses] (stiffness) {\FlowNode{GetPairStiffness}{GetPairStiffness}};
\node[function, below=of stiffness] (velocitydelta) {\FlowNode{ApplySourceVelocityDelta}{ApplySourceVelocityDelta}};
\node[function, below=of velocitydelta] (mass) {\FlowNode{GetParticleMass}{GetParticleMass}};
\node[function, below=of mass] (ledgers) {\FlowNode{UpdateContactLedgers}{UpdateContactLedgers}};
\node[function, below=of ledgers] (setcollisionmom) {\FlowNode{SetCollisionInternalMomentum}{SetCollisionInternalMomentum}};
\node[function, below=of setcollisionmom] (setphase) {\FlowNode{SetCollisionPhase}{SetCollisionPhase}};
\node[function, below=of setphase] (syncmom) {\FlowNode{SyncInternalMomentum}{SyncInternalMomentum}};
\node[function, below=of syncmom] (getmom) {\FlowNode{GetInternalMomentum}{GetInternalMomentum}};
\node[function, below=of getmom] (setmom) {\FlowNode{SetInternalMomentum}{SetInternalMomentum}};

\draw[flow] (contact) -- (begin);
\draw[flow] (begin) -- (init);
\draw[flow] (init) -- (state);
\draw[flow] (state) -- (append);
\draw[flow] (append) -- (slots);
\draw[flow] (slots) -- (phase);
\draw[flow] (phase) -- (collisionmom);
\draw[flow] (collisionmom) -- (geometry);
\draw[flow] (geometry) -- (position);
\draw[flow] (position) -- (startvelocity);
\draw[flow] (startvelocity) -- (reducedmass);
\draw[flow] (reducedmass) -- (parms);

\draw[loop] (parms.west) -- ++(-0.25in,0) |- node[pos=0.16,left,font=\scriptsize]{next contact} (init.west);
\draw[flow] (parms.east) -- ++(0.28in,0) |- node[pos=0.15,right,font=\scriptsize]{all contacts} (weights.south);

\draw[flow] (weights) -- (impulses);
\draw[flow] (impulses) -- (stiffness);
\draw[loop] (stiffness.east) -- ++(0.22in,0) |- (impulses.east);
\draw[flow] (impulses.west) -- ++(-0.22in,0) |- (velocitydelta.west);
\draw[flow] (velocitydelta) -- (mass);
\draw[loop] (mass.east) -- ++(0.22in,0) |- (velocitydelta.east);
\draw[flow] (velocitydelta.west) -- ++(-0.22in,0) |- (ledgers.west);
\draw[flow] (ledgers) -- (setcollisionmom);
\draw[flow] (setcollisionmom) -- (setphase);
\draw[flow] (setphase) -- (syncmom);
\draw[flow] (syncmom) -- (getmom);
\draw[flow] (getmom) -- (setmom);

\end{tikzpicture}
}
\caption{Current ShadowLinDynamics function flow}
\end{figure}

\clearpage
\ShadowFunctionSection{ContactDynamics}

\texttt{ContactDynamics(SourceID)} is the readable top-level dynamics chain for
one source particle.  The source contact list must already contain the
current-frame target IDs.

It resets source accumulators, initializes and calculates each contact, then
performs the source-level stages that require the complete contact set:

\begin{verbatim}
BeginSourceContactParms
for each TargetID:
    InitializeContactState
    CalculateContactParms
CalculateContactWeights
CalculateContactImpulses
ApplySourceVelocityDelta
UpdateContactLedgers
\end{verbatim}

Each stage returns immediately on failure.  The function reads target snapshot
state but writes only source-owned state.

\ShadowFunctionSection{BeginSourceContactParms}

\texttt{BeginSourceContactParms(SourceID)} resets the source-level values that
will be rebuilt from the current contact list:

\begin{itemize}
    \item \texttt{total\_overlap\_area}
    \item \texttt{available\_source\_momentum}
\end{itemize}

These values must be zero before individual contacts contribute their geometry
and closing momentum.

\ShadowFunctionSection{InitializeContactState}

\texttt{InitializeContactState(SourceID, TargetID)} initializes one current
source-target contact.  It obtains the source-owned contact slot from
\texttt{GetContactState}, marks the source as colliding, and updates the
source's reporting target and contact count.  It returns the initialized slot,
or \texttt{None} when initialization fails.

\ShadowFunctionSection{GetContactState}

\texttt{GetContactState(SourceID, TargetID)} creates and populates one
current-frame source-owned contact slot.

The function first calls \texttt{AppendContactSlot} to record the target and
persistent collision state.  It then calls \texttt{GetParticleGeometry} and
stores the contact normal, overlap area, center distance, overlap depth, and
stored internal momentum in the slot.

\ShadowFunctionSection{AppendContactSlot}

\texttt{AppendContactSlot(SourceID, TargetID)} allocates the next available
contact slot owned by the source particle.  The slot records:

\begin{itemize}
    \item target particle ID;
    \item particle-contact type;
    \item persistent collision phase;
    \item active-this-frame state;
    \item previously stored internal momentum.
\end{itemize}

The source's \texttt{contactCount} is incremented only after a slot is filled.
The function returns \texttt{None} when the fixed slot array is full.

\ShadowFunctionSection{GetContactSlots}

\texttt{GetContactSlots(SourceID)} returns the fixed contact-slot array owned by
the source particle.  Newer particle structures use \texttt{contacts}; the
\texttt{gcs} fallback preserves compatibility with the current particle
structure.

\ShadowFunctionSection{GetCollisionPhase}

\texttt{GetCollisionPhase(SourceID, TargetID)} reads the persistent
source-owned phase for one contact.  Collision phases are stored by target ID.
A contact without phase history begins in compression.

\ShadowFunctionSection{GetCollisionInternalMomentum}

\texttt{GetCollisionInternalMomentum(SourceID, TargetID)} reads the persistent
internal momentum owned by one source-target contact ledger.  A missing ledger
or target entry represents zero stored momentum.

\ShadowFunctionSection{GetParticleGeometry}

\texttt{GetParticleGeometry(SourceID, TargetID)} calculates current-frame
contact geometry from frame-start particle positions.

When the two radii overlap, it returns:

\begin{verbatim}
(normal_x, normal_y, normal_z, overlap_area, center_distance)
\end{verbatim}

The normal points from source to target.  The function returns \texttt{None}
when no overlap exists.  Coincident centers use the positive x direction as a
deterministic fallback normal.

\ShadowFunctionSection{GetParticlePosition}

\texttt{GetParticlePosition(ParticleID)} returns the particle's frame-start
position.  It prefers the immutable \texttt{PosLocFrame} snapshot so every
source reads the same frame state.  When no snapshot exists, it reads the
currently active particle position buffer.

\ShadowFunctionSection{CalculateContactParms}

\texttt{CalculateContactParms(SourceID, TargetID, contact\_state)} calculates
one contact's relative normal motion and contribution to source-level totals.

Using frame-start source and target velocities, it calculates:

\begin{equation}
v_{\mathrm{rel},n}
=
(\mathbf{v}_{t}-\mathbf{v}_{s})\cdot\mathbf{n}
\end{equation}

\begin{equation}
p_{\mathrm{closing}}
=
\mu\max(0,-v_{\mathrm{rel},n})
\end{equation}

The relative normal velocity and closing momentum are stored in the contact
slot.  The contact overlap area and closing momentum are also accumulated into
the source totals.

\ShadowFunctionSection{GetStartFrameVelocity}

\texttt{GetStartFrameVelocity(ParticleID)} returns the immutable frame-start
velocity.  It prefers \texttt{VelRadFrame}, preventing one source from observing
another source's velocity write during the same frame.  It falls back to the
particle's current velocity when no snapshot is available.

\ShadowFunctionSection{GetReducedMass}

\texttt{GetReducedMass(SourceID, TargetID)} returns the Newtonian reduced mass:

\begin{equation}
\mu = \frac{m_s m_t}{m_s + m_t}.
\end{equation}

Masses come from the GLSL-shaped \texttt{parms.x} particle field.  An epsilon
floor prevents a zero denominator.

\ShadowFunctionSection{CalculateContactWeights}

\texttt{CalculateContactWeights(SourceID)} distributes the source's available
momentum across its current contacts using source-owned overlap-area weights:

\begin{equation}
w_i = \frac{A_i}{\sum_j A_j}
\end{equation}

\begin{equation}
p_{\mathrm{share},i} = p_{\mathrm{available},s} w_i.
\end{equation}

The source available share is stored on each contact.  This stage does not
inspect target-owned contact state and does not apply an impulse.

\ShadowFunctionSection{CalculateContactImpulses}

\texttt{CalculateContactImpulses(SourceID)} plans compression or release for
each source-owned contact without changing velocity or persistent ledgers.

The geometric raw impulse is:

\begin{equation}
J_{\mathrm{raw}} = q_{\mathrm{pair}} A \Delta t.
\end{equation}

For a closing contact, compression is limited by the contact's weighted source
momentum:

\begin{equation}
J_{\mathrm{compression}}
=
\min(J_{\mathrm{raw}},p_{\mathrm{weighted}}).
\end{equation}

For a non-closing contact, release is limited by stored contact internal
momentum:

\begin{equation}
J_{\mathrm{release}}
=
\min(J_{\mathrm{raw}},p_{\mathrm{internal}}).
\end{equation}

\ShadowFunctionSection{GetPairStiffness}

\texttt{GetPairStiffness(SourceID, TargetID)} returns the nonnegative mean of
the source-owned and target-owned particle stiffness values:

\begin{equation}
q_{\mathrm{pair}} = \max\left(0,\frac{q_s+q_t}{2}\right).
\end{equation}

\ShadowFunctionSection{ApplySourceVelocityDelta}

\texttt{ApplySourceVelocityDelta(SourceID)} converts each planned scalar
contact impulse into a vector along the negative source-to-target normal:

\begin{equation}
\Delta\mathbf{p}_i
=
-(J_{\mathrm{compression},i}+J_{\mathrm{release},i})\mathbf{n}_i.
\end{equation}

All contact momentum deltas are summed before the source velocity is written:

\begin{equation}
\mathbf{v}'_s
=
\mathbf{v}_s+\frac{\sum_i\Delta\mathbf{p}_i}{m_s}.
\end{equation}

Only the source velocity is changed.  The function also records per-contact
momentum deltas and source before/after diagnostics.

\ShadowFunctionSection{GetParticleMass}

\texttt{GetParticleMass(ParticleID)} reads particle mass from
\texttt{parms.x}.  The returned value is bounded below by \texttt{EPSILON}.

\ShadowFunctionSection{UpdateContactLedgers}

\texttt{UpdateContactLedgers(SourceID)} commits planned compression and release
to persistent source-owned contact ledgers:

\begin{equation}
p_{\mathrm{internal},i}'
=
\max(0,
p_{\mathrm{internal},i}
+ J_{\mathrm{compression},i}
- J_{\mathrm{release},i}).
\end{equation}

The function updates each contact's phase and current slot diagnostics, then
rebuilds the source total internal momentum after all contact ledgers are
committed.

\ShadowFunctionSection{SetCollisionInternalMomentum}

\texttt{SetCollisionInternalMomentum(SourceID, TargetID, internal\_momentum)}
stores one source-owned contact internal-momentum value by target ID.  Values at
or below \texttt{EPSILON} are removed from the ledger.

\ShadowFunctionSection{SetCollisionPhase}

\texttt{SetCollisionPhase(SourceID, TargetID, phase)} stores one persistent
source-owned collision phase by target ID.  A drained compression contact is
the default state and therefore requires no stored phase entry.

\ShadowFunctionSection{SyncInternalMomentum}

\texttt{SyncInternalMomentum(SourceID)} rebuilds the source's total internal
momentum from all source-owned contact ledgers and stores the result on the
particle.

\ShadowFunctionSection{GetInternalMomentum}

\texttt{GetInternalMomentum(SourceID)} returns the sum of all stored
source-owned contact internal momentum.  The particle total is used only as a
compatibility fallback when no contact-ledger dictionary exists.

\ShadowFunctionSection{SetInternalMomentum}

\texttt{SetInternalMomentum(SourceID, internal\_momentum)} stores the source
total internal momentum in both the runtime \texttt{internal\_momentum} field
and the GLSL-shaped \texttt{Data.z} field.

\ShadowFunctionSection{Move}

\texttt{Move(SourceID)} integrates one source particle into the inactive
position buffer using its newly updated velocity:

\begin{equation}
\mathbf{x}' = \mathbf{x}+\mathbf{v}\Delta t.
\end{equation}

The function validates the particle ID and time step, updates active-buffer
markers, and reports a lower-bound error when the resulting position leaves the
allowed domain.  The frame-chain owner performs the final position-buffer swap.

\ShadowFunctionSection{isValidParticleID}

\texttt{isValidParticleID(ParticleID)} returns true when the ID is within the
current particle array:

\begin{verbatim}
0 <= ParticleID < len(particles)
\end{verbatim}

\ShadowFunctionSection{SetError}

\texttt{SetError(error\_code)} provides the generic linear-dynamics name for
setting the inherited shared collision error state.  It delegates to the
existing frame-chain error mechanism and returns failure.
